\section{Introduction}

In traditional photography, image acquisition is a direct process: a physical lens focuses light reflected from an object onto a photosensitive sensor, comprised by billions of photosites each of which records the raw intensity values of light in three channels: Red, Green, and Blue. The post-processing needed to then obtain a polished-looking picture are quite straightforward, since the object of interest is the \textit{external} appearance of the subject.

On the other hand, many modern scientific and medical apparatuses require looking inside an object to capture phenomena where direct visual access is not possibile, for example when taking an X-Ray of a patient. In this paradigm, the raw hardware measurements are often highly degraded, indirect, or completely unrecognizable, and the final image must be computationally reconstructed. This constraint gave rise to the interdisciplinary field of \textit{Computational Imaging} (CI), which leverages machine learning and other paradigms to overcome hardware limitations and reconstruct high-quality images from raw, unstructured sensor data. This represents a paradigm shift in visual and diagnostic modalities, transitioning to an integrated approach that couples physical measurement devices with rigorous data-driven and variational algorithms.

An example of an application of computational imaging is MRI undersampling: the idea is to purposely acquire only a fraction of the raw data to speed-up the process of a factor of 2X to 10X, to minimize claustrophobia in patients. The image is then reconstructed from the incomplete data.

Mathematically, the core of Computational Imaging lies in the distinction between the \textit{forward problem} and the \textit{inverse problem}. The forward problem models the underlying physics of the imaging system. Assuming the true, continuous underlying object is represented by $u$, and the discrete, noisy measurements recorded by the hardware are $y$, the acquisition process can be formulated as:
\begin{equation}
    y = Au + \eta
\end{equation}
where $A$ represents the forward operator (e.g., a blurring kernel, a Fourier transform, or a Radon transform) dictated by the system's physics, and $\eta$ denotes noise added by the measurement. The objective of Computational Imaging is to solve the corresponding \textit{inverse problem}: given the corrupted measurements $y$ and the known or estimated forward operator $A$, find an approximation $u'$ that accurately recovers the true state of the object.

Two concrete and complementary archetypes of such inverse problems are explored in this essay. The first is \textit{Image Scanning Microscopy} (ISM), a fluorescence imaging modality in which a two-dimensional detector array replaces the single-element detector of a confocal microscope, yielding a four-dimensional dataset. Here the forward operator is a convolution with a Point Spread Function (PSF), and the reconstruction task consists of recovering a fluorophore map from multiple noisy and shifted views of the same subject. The second is \textit{Computed Tomography} (CT), where a scanner measures the attenuation of X-ray beams passing through the body at various angles, producing raw data known as a sinogram. In this case the forward operator is the Radon transform, and the reconstruction task is to recover a cross-sectional image from a series of indirect projections. Despite their different physical origins and scales, both problems are governed by the same mathematical structure: a linear forward operator whose inversion is ill-posed, requiring principled regularisation or learned priors to yield stable and high-quality reconstructions.

The last half of this essay explores practical approaches to Computed Tomography, implemented using the \textit{DeepInverse} library. First, we will reconstruct a precise internal cross-sections, with the raw data represented as a sonogram, consisting of discrete, noisy line integrals corrupted by noise.  The objective is to formulate a rigorous forward model of the scanning process based on the discrete Radon transform and solve the resulting ill-posed inverse problem to recover the tissue attenuation map. We will implement optimization algorithms for Maximum Likelihood and Maximum a Posteriori estimation.

Then, we will try to learn the prior distribution implicitly from data. We will explore \textit{Plug-and-Play (PnP) Priors}, replacing the proximal operator of the regularizer with a pre-trained deep denoising neural network. This is a hybrid approach that combines known physics with learned image statistics. Lastly, we will train our own denoiser to use it in a PGD network.

\paragraph{Sources}

The main sources for this essay are the following:

\begin{itemize}
    \item Prof. Luca Calatroni's mini-course slides "Computational imaging \& learning: from physics to data-driven methods";
    \item C. A. Bouman, Foundations of Computational Imaging: A Model-Based Approach, SIAM, 2022 \cite{bouman_foundations_2022};
    \item G. Aubert, P. Kornprobst, Mathematical Problems in Image Processing: PDEs and the COV, II edition, Springer, 2000 \cite{aubert_mathematical_2008};
    \item Tony F. Chan, J. Shen, Image Processing and Analysis: variational, PDE, wavelet and stochastic methods, SIAM, 2005 \cite{chan_image_2005};
    \item Tristan van Leeuwen's blog \cite{van_leeuwen_1_nodate};
    \item Johannes Meyer's blog \cite{meyer_computational_nodate};
    \item Numerical tours by G. Peyré \cite{peyre_numerical_nodate};
    \item DeepInverse tutorials \cite{deepinverse_examples_nodate}.
\end{itemize}